\newpage
\appendix
\section*{Anexos}
\addcontentsline{toc}{section}{Anexos}
\renewcommand{\thesubsection}{\Alph{subsection}}
\setcounter{subsection}{0}

% =============================================================================
\subsection{Matriz de revisión bibliográfica}
\label{anexo:matriz}

La tabla~\ref{tab:matriz-biblio} sistematiza las fuentes utilizadas. Todas
fueron verificadas individualmente contra su fuente primaria en septiembre de
2026; se consigna el identificador digital de objeto de cada una para permitir
la verificación independiente por parte del evaluador.

\begin{table}[H]
\centering
\caption{Matriz de revisión bibliográfica.}
\label{tab:matriz-biblio}
\scriptsize
\begin{tabular}{@{}p{2.4cm}p{1.5cm}p{3.1cm}p{3.4cm}p{3.3cm}@{}}
\toprule
\textbf{Fuente} & \textbf{Eje} & \textbf{Aporte principal} &
\textbf{Relación con el proyecto} & \textbf{DOI}\\
\midrule
\textcite{kalman1960} & Estimación &
Estimador recursivo de varianza mínima para sistemas lineales &
Fundamento algorítmico del módulo (OE3) &
10.1115/1.3662552\\
\addlinespace
\textcite{simon2006} & Estimación &
Tratamiento moderno de la estimación óptima y convergencia de Riccati &
Sustento de la convergencia asintótica (OE3, OE4) &
10.1002/0470045345\\
\addlinespace
\textcite{lee2008} & Gemelos &
Formulación del paradigma de sistemas ciberfísicos &
Encuadre del acoplamiento físico-computacional &
10.1109/ISORC.2008.25\\
\addlinespace
\textcite{uhlemann2017} & Gemelos &
Gemelo digital como realización del sistema ciberfísico &
Contrapunto: sitúa la brecha en el \emph{middleware} &
10.1016/j.procir.2016.11.152\\
\addlinespace
\textcite{tao2019} & Gemelos &
Gemelo digital para pronóstico y gestión de salud de equipos &
Aplicación industrial de referencia; asume telemetría limpia &
10.1016/j.cirp.2019.04.049\\
\addlinespace
\textcite{fuller2020} & Gemelos &
Tecnologías habilitantes y desafíos abiertos &
Reconoce la calidad de datos como problema no resuelto &
10.1109/ACCESS.2020.2998358\\
\addlinespace
\textcite{liu2021} & Gemelos &
Revisión de conceptos, tecnologías y aplicaciones industriales &
Estado de adopción industrial del paradigma &
10.1016/j.jmsy.2020.06.017\\
\addlinespace
\textcite{sepasgozar2021} & Gemelos &
Distinción entre sombra digital y gemelo digital &
Criterio conceptual central del proyecto &
10.3390/buildings11040151\\
\addlinespace
\textcite{apelt2002} & Minería &
Medición inferencial de parámetros de molino SAG por estimación de estados &
Precedente directo de reconstrucción de estados en conminución &
10.1016/S0892-6875(02)00229-7\\
\addlinespace
\textcite{qu2023} & Minería &
Revisión comprehensiva de gemelos digitales en la industria de minerales &
Evidencia de adopción fragmentada; sitúa el vacío &
10.1080/25726668.2023.2257479\\
\addlinespace
\textcite{barzaga2025} & Minería &
Filtros de Kalman extendido y \emph{unscented} en lixiviación de cobre &
Precedente en minería del cobre; contraste lineal / no lineal &
10.3390/min15050521\\
\addlinespace
\textcite{ebad2026} & Minería &
Revisión sistemática de 68 estudios de gemelos digitales con IA en minería &
Documenta la escasa validación industrial: sustenta el vacío &
10.3390/technologies14050269\\
\bottomrule
\end{tabular}
\end{table}

\subsubsection*{Nota sobre depuración bibliográfica}

Durante la consolidación de esta entrega se verificó la totalidad de las
referencias empleadas en las fases previas. Cuatro de ellas no pudieron ser
localizadas en sus revistas de origen ni en repositorios bibliográficos, por lo
que fueron retiradas del trabajo y sustituidas por fuentes verificadas de
pertinencia equivalente o superior. Adicionalmente, se corrigió la atribución
editorial de \textcite{sepasgozar2021}, que en entregas anteriores aparecía
asignada a una revista distinta de aquella en que fue efectivamente publicada.
Esta depuración se documenta aquí en cumplimiento del principio de integridad
académica declarado en la sección~\ref{sec:etica}.

% =============================================================================
\subsection{Matriz de correspondencia entre el informe y el artículo científico}
\label{anexo:correspondencia}

\begin{table}[H]
\centering
\caption{Correspondencia entre secciones del informe final y del artículo
científico en formato IEEE.}
\label{tab:correspondencia}
\footnotesize
\begin{tabular}{@{}p{4.6cm}p{4.6cm}p{4.2cm}@{}}
\toprule
\textbf{Sección del informe} & \textbf{Sección del artículo} &
\textbf{Estado de trazabilidad}\\
\midrule
\S\ref{sec:fenomeno}--\ref{sec:pregunta} Fenómeno, problema y pregunta &
I. Introduction & Consolidado\\
\addlinespace
\S\ref{sec:estado-arte} Estado del arte y vacío &
II. Related Work & Consolidado; misma bibliografía verificada\\
\addlinespace
\S\ref{sec:unidad} Unidad de análisis &
III. Problem Formulation & Modelo y matrices unificados\\
\addlinespace
\S\ref{sec:enfoque}--\ref{sec:procedimientos} Metodología &
IV. Methodology & Consolidado\\
\addlinespace
\S\ref{sec:analisis} Estrategia de análisis &
IV.D. Evaluation Metrics and Tests & Métricas y pruebas unificadas\\
\addlinespace
\S\ref{sec:resultados} Resultados &
V. Results & Mismas figuras y valores numéricos\\
\addlinespace
\S\ref{sec:discusion} Discusión &
VI. Discussion & Consolidado\\
\addlinespace
\S\ref{sec:conclusiones} Conclusiones &
VII. Conclusion & Consolidado\\
\bottomrule
\end{tabular}
\end{table}

La trazabilidad entre ambos productos es verificable y no meramente declarada:
comparten el mismo archivo de referencias, las mismas figuras generadas por el
código anexo y los mismos valores numéricos, de modo que cualquier discrepancia
sería detectable por inspección directa.

% =============================================================================
\subsection{Ajustes derivados de la retroalimentación docente}
\label{anexo:ajustes}

\begin{table}[H]
\centering
\caption{Registro de ajustes incorporados en la entrega final.}
\label{tab:ajustes}
\footnotesize
\begin{tabular}{@{}p{1.8cm}p{4.6cm}p{4.6cm}p{2.4cm}@{}}
\toprule
\textbf{Origen} & \textbf{Observación recibida} & \textbf{Ajuste realizado} &
\textbf{Sección}\\
\midrule
Fases 1--3 &
Lista de referencias en formato numérico cuando la pauta exige APA 7
autor-fecha &
Unificación completa bajo APA 7 mediante gestor bibliográfico &
Referencias\\
\addlinespace
Fase 1 &
Ausencia de matriz de revisión bibliográfica &
Incorporación de la matriz con caracterización y DOI de cada fuente &
Anexo~\ref{anexo:matriz}\\
\addlinespace
Fase 1 &
El plan del proyecto era en realidad una sección de alcances; faltaban
resultados esperados, hitos, avances ABP y trazabilidad al artículo &
Formulación completa del plan con etapas, entregables, indicadores, hitos,
criterios de seguimiento y viabilidad &
\S\ref{sec:plan}\\
\addlinespace
Fase 1 &
Cita de un \emph{working paper} propio no publicado como respaldo externo &
Retirada; sustituida por fuentes indexadas verificadas &
Anexo~\ref{anexo:matriz}\\
\addlinespace
Fase 2 &
Trazabilidad con el artículo declarada pero no verificable &
Matriz de correspondencia explícita y recursos compartidos entre ambos
productos &
Anexo~\ref{anexo:correspondencia}\\
\addlinespace
Fase 2 &
Enunciar los objetivos específicos y amarrarlos a cada eje conceptual del marco
teórico &
Cada eje del marco teórico declara explícitamente el objetivo que sustenta &
\S\ref{sec:marco}\\
\addlinespace
Fase 3 &
Consideraciones éticas sin estructura riesgo--mitigación y sin nombrar el riesgo
central del estudio &
Reestructuración en cinco riesgos con su acción de mitigación, encabezados por
la sobreatribución de madurez tecnológica &
\S\ref{sec:etica}\\
\addlinespace
Fase 3 &
Revisar la unidad de la varianza $R$ &
Declaración explícita de unidades de todos los parámetros con su fundamento
físico &
Tabla~\ref{tab:parametros}\\
\addlinespace
Fase 3 &
Verificar el supuesto de normalidad de la prueba t &
Prueba de Shapiro-Wilk aplicada; ante su rechazo se incorporó el contraste no
paramétrico de Wilcoxon &
\S\ref{sec:analisis}, \S\ref{sec:resultados}\\
\addlinespace
Fase 3 &
Incoherencia dimensional: MSE en \si{\metre\tothe{6}} e intervalo de confianza
en \si{\metre\cubed} &
Separación explícita de métricas; toda la inferencia se funda sobre el RMSE en
\si{\metre\cubed} &
\S\ref{sec:analisis}\\
\addlinespace
Fase 3 &
Unificación del título del artículo científico &
Título del artículo alineado con el foco en chancado primario &
Artículo IEEE\\
\addlinespace
Fase 3 &
Análisis de sensibilidad declarado en los criterios de calidad pero no
ejecutado &
Implementación del barrido de degradación y de la malla de desajuste
$\mathbf{Q}/\mathbf{R}$ &
\S\ref{sec:resultados}\\
\bottomrule
\end{tabular}
\end{table}

Se incorporaron además dos ajustes que no provienen de la retroalimentación sino
de la implementación efectiva del modelo: la recalibración de la covarianza de
proceso $Q_{22}$ por implausibilidad física
(\S\ref{sec:recalibracion}) y la unificación de la condición inicial del
sistema, que presentaba valores distintos en documentos previos.

% =============================================================================
\subsection{Anexo digital: código fuente y registros de trazabilidad}
\label{anexo:codigo}

Se adjunta el paquete completo del módulo, cuya estructura es la siguiente:

\begin{table}[H]
\centering
\caption{Contenido del anexo digital.}
\label{tab:anexo-digital}
\footnotesize
\begin{tabular}{@{}p{4.6cm}p{8.6cm}@{}}
\toprule
\textbf{Componente} & \textbf{Descripción}\\
\midrule
\texttt{nucleo\_kalman.py} &
Núcleo del estimador, genérico y desacoplado del proceso físico\\
\texttt{modelo\_tolva.py} &
Modelo de referencia de la tolva, parametrización y degradación SCADA\\
\texttt{analisis.py} &
Métricas de error, simulación de Monte Carlo e inferencia estadística\\
\texttt{sensibilidad.py} &
Barrido de degradación y malla de desajuste de covarianzas\\
\texttt{figuras.py} &
Generación de las figuras del informe y del artículo\\
\texttt{main.py} &
Punto de entrada único del experimento\\
\texttt{README.md} &
Documentación, parametrización y declaración de madurez tecnológica\\
\texttt{resultados/logs/} &
Cuatro registros CSV de trazabilidad\\
\texttt{resultados/figuras/} &
Seis figuras en formato vectorial y de mapa de bits\\
\texttt{resultados/resultados.json} &
Parametrización, entorno de ejecución y resumen consolidado\\
\bottomrule
\end{tabular}
\end{table}

La ejecución del archivo \texttt{main.py} en un entorno con Python~3, NumPy,
SciPy y Matplotlib reproduce la totalidad de los resultados reportados en este
informe.
